% =============================================================================
% CAPÍTULO 3: METODOLOGIA
% =============================================================================

\chapter{Metodologia}

\section{Visão Geral da Metodologia}

A metodologia proposta neste trabalho é desenvolvida em quatro etapas principais, conforme ilustrado na Figura \ref{fig:metodologia_geral}:

\begin{enumerate}
    \item \textbf{Modelagem e Simulação:} Desenvolvimento de modelos numéricos e geração de dados simulados.
    \item \textbf{Análise HOS:} Aplicação de técnicas de espectros de ordem superior aos dados simulados.
    \item \textbf{Validação Física:} Comparação dos resultados com literatura científica.
    \item \textbf{Classificação:} Desenvolvimento de classificadores de aprendizado de máquina.
\end{enumerate}

\begin{figure}[H]
    \centering
    \includegraphics[width=0.8\textwidth]{figures/metodologia_geral.png}
    \caption{Fluxograma da metodologia proposta}
    \label{fig:metodologia_geral}
\end{figure}

\section{Ambiente Computacional}

\subsection{Ferramentas Utilizadas}

O desenvolvimento da metodologia utiliza as seguintes ferramentas computacionais:

\begin{itemize}
    \item \textbf{ROSS (Rotordynamic Open-Source Software):} Biblioteca Python para simulação de sistemas rotativos.
    \item \textbf{Python:} Linguagem de programação principal para implementação dos algoritmos.
    \item \textbf{Bibliotecas científicas:} NumPy, SciPy, Matplotlib, Pandas para análise de dados.
    \item \textbf{Scikit-learn:} Para implementação de algoritmos de aprendizado de máquina.
\end{itemize}

\subsection{Configuração do Ambiente}

O ambiente de desenvolvimento foi configurado com as seguintes especificações:

\begin{lstlisting}[language=Python, caption=Configuração do ambiente Python]
# Requisitos do sistema
python >= 3.8
ross >= 0.2.0
numpy >= 1.21.0
scipy >= 1.7.0
matplotlib >= 3.4.0
pandas >= 1.3.0
scikit-learn >= 1.0.0
\end{lstlisting}

\section{Modelagem Numérica}

\subsection{Geometria do Sistema}

O sistema rotativo modelado consiste em um eixo de aço com as seguintes características:

\begin{table}[H]
    \centering
    \caption{Propriedades do material do eixo}
    \begin{tabular}{lc}
        \toprule
        Propriedade & Valor \\
        \midrule
        Módulo de elasticidade & 200 GPa \\
        Densidade & 7850 kg/m³ \\
        Coeficiente de Poisson & 0.3 \\
        \bottomrule
    \end{tabular}
    \label{tab:propriedades_material}
\end{table}

\subsection{Modelagem de Falhas}

\subsubsection{Modelo de Trinca}

O modelo de trinca implementado considera o efeito de "respiração", onde a rigidez da trinca varia com a rotação do eixo. A rigidez efetiva é modelada como:

\begin{equation}
k_{crack}(\theta) = k_0 \left[ 1 - \alpha \cos(\theta) \right]
\end{equation}

onde $k_0$ é a rigidez nominal, $\alpha$ é o parâmetro de severidade da trinca e $\theta$ é o ângulo de rotação.

\subsubsection{Modelo de Desalinhamento}

O desalinhamento é modelado através de forças desbalanceadas aplicadas nos mancais:

\begin{equation}
F_{misalignment} = m \omega^2 r \sin(\omega t + \phi)
\end{equation}

onde $m$ é a massa desbalanceada, $\omega$ é a velocidade angular, $r$ é o raio de desbalanceamento e $\phi$ é o ângulo de fase.

\section{Simulações Numéricas}

\subsection{Parâmetros de Simulação}

As simulações foram realizadas com os seguintes parâmetros:

\begin{itemize}
    \item \textbf{Velocidades de rotação:} 1000, 2000, 3000, 4000, 5000 RPM
    \item \textbf{Severidade das falhas:} 5\%, 10\%, 15\%, 20\% da rigidez nominal
    \item \textbf{Tempo de simulação:} 10 segundos
    \item \textbf{Frequência de amostragem:} 2048 Hz
\end{itemize}

\subsection{Processamento de Dados}

Os dados de vibração simulados são processados através das seguintes etapas:

\begin{enumerate}
    \item \textbf{Pré-processamento:} Remoção de tendências e filtragem.
    \item \textbf{Segmentação:} Divisão dos sinais em janelas de análise.
    \item \textbf{Análise espectral:} Cálculo de espectros de potência.
    \item \textbf{Análise HOS:} Cálculo de bi-espectros e tri-espectros.
\end{enumerate}

\section{Análise de Espectros de Ordem Superior}

\subsection{Implementação do Bi-espectro}

O bi-espectro é calculado utilizando a seguinte implementação:

\begin{lstlisting}[language=Python, caption=Cálculo do bi-espectro]
def calculate_bispectrum(signal, fs, nperseg=1024):
    """
    Calcula o bi-espectro de um sinal
    
    Parameters:
    signal: array-like, sinal de entrada
    fs: float, frequência de amostragem
    nperseg: int, número de pontos por segmento
    
    Returns:
    f: array, frequências
    bispectrum: array, bi-espectro
    """
    from scipy.signal import welch
    from scipy.fft import fft
    
    # Segmentação do sinal
    segments = []
    for i in range(0, len(signal) - nperseg, nperseg // 2):
        segments.append(signal[i:i + nperseg])
    
    # Cálculo do bi-espectro para cada segmento
    bispectra = []
    for segment in segments:
        X = fft(segment)
        bispectrum = np.outer(X, X) * np.conj(X)
        bispectra.append(bispectrum)
    
    # Média dos bi-espectros
    bispectrum_mean = np.mean(bispectra, axis=0)
    
    return bispectrum_mean
\end{lstlisting}

\subsection{Extração de Características}

As características extraídas dos bi-espectros incluem:

\begin{itemize}
    \item \textbf{Magnitude máxima:} Valor máximo do bi-espectro.
    \item \textbf{Frequências de pico:} Coordenadas dos picos principais.
    \item \textbf{Energia total:} Soma dos valores absolutos do bi-espectro.
    \item \textbf{Entropia:} Medida de complexidade do espectro.
\end{itemize}

\section{Validação Física}

\subsection{Estratégia de Validação}

A validação física é realizada através de:

\begin{enumerate}
    \item \textbf{Comparação com literatura:} Análise de padrões espectrais reportados em estudos experimentais.
    \item \textbf{Análise de sensibilidade:} Verificação da resposta do sistema a variações nos parâmetros.
    \item \textbf{Verificação de consistência:} Confirmação de que os resultados seguem princípios físicos conhecidos.
\end{enumerate}

\subsection{Métricas de Validação}

As seguintes métricas são utilizadas para validação:

\begin{itemize}
    \item \textbf{Correlação espectral:} Comparação de padrões de frequência.
    \item \textbf{Erro relativo:} Diferença percentual entre valores simulados e experimentais.
    \item \textbf{Análise estatística:} Testes de significância para validação quantitativa.
\end{itemize}

\section{Classificação por Aprendizado de Máquina}

\subsection{Preparação dos Dados}

Os dados são preparados para classificação através de:

\begin{enumerate}
    \item \textbf{Normalização:} Padronização das características extraídas.
    \item \textbf{Divisão treino-teste:} Separação dos dados em conjuntos de treinamento e teste.
    \item \textbf{Validação cruzada:} Implementação de k-fold cross-validation.
\end{enumerate}

\subsection{Algoritmos de Classificação}

Os seguintes algoritmos são implementados e comparados:

\begin{itemize}
    \item \textbf{Support Vector Machine (SVM):} Com kernel RBF.
    \item \textbf{Random Forest:} Ensemble de árvores de decisão.
    \item \textbf{Neural Networks:} Rede neural multicamadas.
    \item \textbf{Gradient Boosting:} XGBoost para classificação.
\end{itemize}

\subsection{Métricas de Avaliação}

O desempenho dos classificadores é avaliado através de:

\begin{itemize}
    \item \textbf{Acurácia:} Proporção de classificações corretas.
    \item \textbf{Precisão:} Proporção de verdadeiros positivos.
    \item \textbf{Recall:} Proporção de casos positivos identificados.
    \item \textbf{F1-Score:} Média harmônica entre precisão e recall.
    \item \textbf{Matriz de confusão:} Visualização das classificações.
\end{itemize}

\section{Plano de Experimentos}

\subsection{Experimentos Planejados}

O plano de experimentos inclui:

\begin{enumerate}
    \item \textbf{Experimento 1:} Validação do modelo numérico com dados da literatura.
    \item \textbf{Experimento 2:} Análise de sensibilidade dos parâmetros de falha.
    \item \textbf{Experimento 3:} Comparação de diferentes algoritmos HOS.
    \item \textbf{Experimento 4:} Avaliação de classificadores de aprendizado de máquina.
    \item \textbf{Experimento 5:} Análise de robustez com ruído adicionado.
\end{enumerate}

\subsection{Análise Estatística}

A análise estatística dos resultados inclui:

\begin{itemize}
    \item \textbf{Testes de normalidade:} Shapiro-Wilk para verificar distribuição normal.
    \item \textbf{Análise de variância:} ANOVA para comparar grupos.
    \item \textbf{Testes post-hoc:} Tukey HSD para comparações múltiplas.
    \item \textbf{Análise de correlação:} Pearson e Spearman para relações entre variáveis.
\end{itemize}
